
complexes de Koszul associés. Le complexe de Koszul associé à la suite \(f_1,\ldots,f_p,g_1,\ldots,g_r\) est isomorphe au complexe produit tensoriel \(\underline{K}_\bullet(u)\otimes_{\underline{O}_X}\underline{K}_\bullet(v)\). Comme c'est une résolution de \(\underline{O}_Y\) par des \(\underline{O}_X\)-Modules libres de type fini, les \(\underline{\Tor}_i^{\underline{O}_X}(\underline{O}_Y,\underline{O}_Z)\) sont les faisceaux d'homologie du complexe \(\underline{K}_\bullet(u)\otimes_{\underline{O}_X}\underline{K}_\bullet(v)\otimes_{\underline{O}_X}\underline{O}_Z\). On peut le considérer comme le produit tensoriel des complexes \(\underline{K}_\bullet(u)\) et \(\underline{K}_\bullet(v)\otimes_{\underline{O}_X}\underline{O}_Z\), et utiliser pour calculer son homologie les suites spectrales associées à un bicomplexe (EGA \(0_{\mathrm{III}}\) (11.3.5)), qui ont comme aboutissement l'homologie du complexe simple associé.

On considère la suite spectrale telle que \({}'E^2_{pq}=H^I_p\bigl(H^{II}_q(K_{\bullet\bullet})\bigr)\), avec les notations de EGA \(0_{\mathrm{III}}\) (11.3.5), le premier indice étant relatif au complexe \(\underline{K}_\bullet(u)\) et le second au complexe \(\underline{K}_\bullet(v)\otimes_{\underline{O}_X}\underline{O}_Z\). Or le complexe \(\underline{K}_\bullet(v)\otimes_{\underline{O}_X}\underline{O}_Z\) est le complexe de Koszul sur \(\underline{O}_Z\) associé à la suite de sections \(\overline g_1,\ldots,\overline g_r\), où les \(\overline g_i\) sont les images des \(g_i\) dans \(\underline{O}_Z\), et ce complexe est une résolution de \(\underline{O}_T\) sur \(\underline{O}_Z\) car les immersions \(Z\to X\), et \(T\to Z\to X\) sont régulières (cf. démonstration de 1.7 ii). Comme \(\bigwedge^p\underline{E}\) est libre sur \(\underline{O}_X\), le complexe \(\bigwedge^p\underline{E}\otimes_{\underline{O}_X}(\underline{K}_\bullet(v)\otimes_{\underline{O}_X}\underline{O}_Z)\) est une résolution de \(\bigwedge^p\underline{E}\otimes_{\underline{O}_X}\underline{O}_T\). Par conséquent la suite spectrale dégénère, et les \(\underline{\Tor}_i^{\underline{O}_X}(\underline{O}_Y,\underline{O}_Z)\) sont les faisceaux d'homologie du complexe \(\underline{K}_\bullet(u)\otimes_{\underline{O}_X}\underline{O}_T\). Or la différentielle de \(\underline{K}_\bullet(u)\) s'exprime en termes des \(f_i\), qui annulent \(\underline{O}_T\); par suite, la différentielle de \(\underline{K}_\bullet(u)\otimes_{\underline{O}_X}\underline{O}_T\) est nulle, et on obtient donc les isomorphismes:
\begin{equation}\tag{2.5.2}
\underline{\Tor}_i^{\underline{O}_X}(\underline{O}_Y,\underline{O}_Z)
\xrightarrow{\sim}
\underline{O}_T\otimes_{\underline{O}_X}\bigwedge^i\underline{E}
\longrightarrow
\bigwedge^i(\underline{O}_T)^p.
\end{equation}

En particulier, on trouve pour \(i=1\):
\[
\underline{\Tor}_1^{\underline{O}_X}(\underline{O}_Y,\underline{O}_Z)\xrightarrow{\sim}(\underline{O}_T)^p,
\]
d'où pour tout \(i\):
\begin{equation}\tag{2.5.3}
\underline{\Tor}_i^{\underline{O}_X}(\underline{O}_Y,\underline{O}_Z)
\xrightarrow{\sim}
\bigwedge^i\underline{\Tor}_1^{\underline{O}_X}(\underline{O}_Y,\underline{O}_Z).
\end{equation}

